\documentclass[12pt]{article}
\usepackage{fullpage,amsmath,amssymb,amsthm,hyperref,amsfonts}
\usepackage[T1]{fontenc}
\usepackage[utf8]{inputenc}

%% Theorem environments
\newtheorem{theorem}{Theorem}
\newtheorem{lemma}[theorem]{Lemma}
\newtheorem{proposition}[theorem]{Proposition}
\newtheorem{corollary}[theorem]{Corollary}
\theoremstyle{definition}
\newtheorem{definition}[theorem]{Definition}
\newtheorem{example}[theorem]{Example}
\newtheorem{remark}[theorem]{Remark}

%% Notation macros
\newcommand{\calO}{\mathcal{O}}
\newcommand{\Sp}{\mathrm{Sp}}
\newcommand{\SpG}{\mathrm{Sp}^G}
\newcommand{\Sub}{\mathrm{Sub}}
\newcommand{\rhoH}{\rho_H}
\newcommand{\phiH}{\Phi^H}
\newcommand{\fF}{\mathcal{F}}
\newcommand{\floor}[1]{\lfloor #1 \rfloor}

\title{Slice Filtration for $N_\infty$ Operads:\\
       A Characterization via Geometric Fixed Points}
\author{}
\date{April 29, 2026}

\begin{document}
\maketitle

\begin{abstract}
Fix a finite group $G$. Given an $N_\infty$ operad $\calO$ with transfer system
$\mathcal{T}_\calO$, we define the $\calO$-adapted slice filtration on the $G$-equivariant
stable $\infty$-category $\SpG$ and prove: a connective $G$-spectrum $X$ is
$\calO$-$n$-slice-connected if and only if for every $\calO$-admissible subgroup $H \leq G$,
the geometric fixed-point spectrum $\Phi^H(X)$ is $(\floor{n/|H|} - 1)$-connected.
\end{abstract}

\tableofcontents

\section{Background and Definitions}

Throughout, $G$ denotes a fixed finite group. We work in the $G$-equivariant stable
$\infty$-category $\SpG$ in the sense of \cite{HHR2016,LMS1986}; all spectra are
$G$-equivariant unless stated otherwise.

\subsection{Transfer Systems}

Write $\Sub(G)$ for the lattice of subgroups of $G$ ordered by inclusion.

\begin{definition}[Transfer system {\cite{BlumbergHill2015}}]
\label{def:transfer-system}
A \emph{$G$-transfer system} is a collection
$\mathcal{T} \subseteq \{(H, K) \mid H \leq K \leq G\}$
satisfying:
\begin{enumerate}
\item[(T1)] \textbf{Identity:} $(H, H) \in \mathcal{T}$ for every $H \leq G$.
\item[(T2)] \textbf{Conjugation invariance:} If $(H, K) \in \mathcal{T}$ and $g \in G$,
  then $(gHg^{-1}, gKg^{-1}) \in \mathcal{T}$.
\item[(T3)] \textbf{Restriction:} If $(H, K) \in \mathcal{T}$ and $L \leq K$, then
  $(L \cap H, L) \in \mathcal{T}$.
\end{enumerate}
We write $H \leq_{\mathcal{T}} K$ to mean $(H,K) \in \mathcal{T}$.
\end{definition}

\begin{remark}\label{rem:extremal-ts}
The minimal transfer system $\mathcal{T}_{\min} = \{(H,H) \mid H \leq G\}$ and the
maximal transfer system $\mathcal{T}_{\max} = \{(H,K) \mid H \leq K \leq G\}$ correspond,
respectively, to the trivial operad (no norms beyond identity) and to an $E_\infty$-operad
(all norms available).
\end{remark}

\subsection{\texorpdfstring{$N_\infty$}{N-infinity} Operads}

\begin{definition}[$N_\infty$ operad {\cite{BlumbergHill2015}}]
\label{def:ninfty-operad}
An \emph{$N_\infty$ operad} is a $G$-equivariant operad $\calO$ in the category of
$G$-spaces such that for every subgroup $H \leq G$ the $H$-fixed spaces $\calO(n)^H$ are
contractible ($\calO$ is $\Sigma$-free and genuine), and such that the collection of
$\calO$-admissible $H$-sets (those $H$-sets $T$ for which the structure map $\calO(T) \to
\calO(|T|)$ is $H$-equivariantly contractible) is closed under conjugation, restriction, and
disjoint union.

The \emph{transfer system of $\calO$} is
\[
  \mathcal{T}_\calO := \bigl\{(H,K) \bigm| K/H \text{ is an admissible $H$-set for }
  \calO\bigr\},
\]
where $K/H$ carries the $H$-action by left multiplication. The pair $(H,K)$ belongs to
$\mathcal{T}_\calO$ if and only if the norm map $N_H^K$ is available in $\calO$-algebra
structures.
\end{definition}

The classification theorem of Blumberg--Hill \cite{BlumbergHill2015} (proved in final form
in \cite{Rubin2021}) states that the homotopy type of an $N_\infty$ operad is completely
determined by its transfer system, and every $G$-transfer system is realized by some
$N_\infty$ operad.

\subsection{\texorpdfstring{$\calO$}{O}-Admissible Subgroups}

\begin{definition}[$\calO$-admissible subgroup]
\label{def:admissible}
A subgroup $H \leq G$ is \emph{$\calO$-admissible} if $(e, H) \in \mathcal{T}_\calO$, where
$e = \{1\}$ is the trivial subgroup. Equivalently, the norm map $N_e^H$ (the $H$-fold norm
from spectra to $G$-spectra) is present in $\calO$-algebra structures.
\end{definition}

\begin{remark}
Every subgroup $e$ is always $\calO$-admissible by axiom (T1) since $(e,e) \in
\mathcal{T}_\calO$. The group $G$ itself is $\calO$-admissible if and only if the full norm
$N_e^G$ is present in $\calO$.
\end{remark}

\begin{lemma}[Restriction of admissibility]
\label{lem:restriction-admissibility}
Let $H$ be $\calO$-admissible and $K \leq H$. Then $K$ is also $\calO$-admissible.
\end{lemma}

\begin{proof}
Since $(e, H) \in \mathcal{T}_\calO$, applying axiom (T3) with $L = K \leq H$ gives
$(K \cap e, K) = (e, K) \in \mathcal{T}_\calO$.
\end{proof}

\subsection{Regular Representation and Slice Cells}

\begin{definition}[Regular representation]
\label{def:regular-rep}
For a subgroup $H \leq G$, the \emph{regular representation} $\rho_H$ is the real
$H$-representation on $\mathbb{R}[H]$ (the free $\mathbb{R}$-module on the elements of $H$),
where $H$ acts by left multiplication. We have $\dim_\mathbb{R}(\rho_H) = |H|$. The
$H$-fixed subspace is
\[
  \rho_H^H = \mathbb{R}\cdot\!\Bigl(\sum_{h\in H} h\Bigr) \cong \mathbb{R},
\]
a one-dimensional trivial representation, so $\dim_\mathbb{R}(\rho_H^H) = 1$.
\end{definition}

\begin{lemma}[Restriction of regular representations]
\label{lem:restriction-regular}
If $K \leq H$, then $\mathrm{Res}^H_K \rho_H \cong |H/K| \cdot \rho_K$ as $K$-representations.
\end{lemma}

\begin{proof}
Write $H = \bigsqcup_{i=1}^{|H/K|} K h_i$ as a disjoint union of left $K$-cosets. Then
\[
  \mathbb{R}[H] = \bigoplus_{i=1}^{|H/K|} \mathbb{R}[K h_i],
\]
and each summand $\mathbb{R}[K h_i]$ is isomorphic to $\mathbb{R}[K]$ as a $K$-module via
$k h_i \mapsto k$. (This map is $K$-equivariant: for $k' \in K$,
$k'(k h_i) = (k'k) h_i \mapsto k'k$, and $k' \cdot k = k'k$ in $\mathbb{R}[K]$.)
Hence $\mathrm{Res}^H_K \rho_H \cong |H/K| \cdot \rho_K$.
\end{proof}

\begin{definition}[$\calO$-slice cells]
\label{def:O-slice-cells}
For an $\calO$-admissible subgroup $H \leq G$ and $k \geq 0$, the \emph{$\calO$-slice cell
of dimension $k|H|$} is
\[
  C^{\calO}(H, k) := G_+ \wedge_H S^{k\rho_H},
\]
the $G$-spectrum obtained by inducing the $H$-suspension spectrum of the sphere $S^{k\rho_H}$
from $H$ to $G$.
\end{definition}

\subsection{Geometric Fixed Points}

\begin{definition}[Geometric fixed points {\cite{LMS1986,HHR2016}}]
\label{def:geom-fixed}
For $H \leq G$, let $\mathcal{F}[H]$ be the family of \emph{proper} subgroups of $H$.
The classifying $H$-space $E\mathcal{F}[H]$ is the $H$-CW complex (unique up to
$H$-equivariant homotopy) with
\[
  E\mathcal{F}[H]^K \simeq \begin{cases} * & K \in \mathcal{F}[H], \\ \emptyset & K = H. \end{cases}
\]
Its unreduced suspension $\widetilde{E}\mathcal{F}[H]$ is the cofiber of
$E\mathcal{F}[H]_+ \to S^0$; it satisfies
\[
  \widetilde{E}\mathcal{F}[H]^K \simeq \begin{cases} * & K \in \mathcal{F}[H], \\
  S^0 & K = H. \end{cases}
\]
The \emph{geometric fixed-point functor} $\Phi^H \colon \SpG \to \Sp$ is
\begin{equation}
\label{eq:geom-fixed-def}
  \Phi^H(X) := \bigl(\widetilde{E}\mathcal{F}[H] \wedge X\bigr)^H.
\end{equation}
\end{definition}

\begin{lemma}[Geometric fixed points of representation spheres]
\label{lem:geom-fixed-sphere}
For any finite-dimensional real $H$-representation $V$,
\begin{equation}
\label{eq:geom-fixed-sphere}
  \Phi^H(S^V) \simeq S^{V^H},
\end{equation}
where $V^H \subseteq V$ is the subspace of $H$-fixed vectors.
In particular,
\begin{equation}
\label{eq:geom-fixed-regular}
  \Phi^H(S^{k\rho_H}) \simeq S^k
  \quad \text{for every } k \geq 0,
\end{equation}
since $\dim_\mathbb{R}(\rho_H^H) = 1$.
\end{lemma}

\begin{proof}
We apply the cofiber sequence $E\mathcal{F}[H]_+ \to S^0 \to \widetilde{E}\mathcal{F}[H]$
and smash with $S^V$:
\begin{equation}
\label{eq:cofiber-smash}
  E\mathcal{F}[H]_+ \wedge S^V \longrightarrow S^V
  \longrightarrow \widetilde{E}\mathcal{F}[H] \wedge S^V.
\end{equation}
Applying $(-)^H$ to \eqref{eq:cofiber-smash}:

\textbf{Step 1.} $(E\mathcal{F}[H]_+ \wedge S^V)^H$: since $(E\mathcal{F}[H])^H = \emptyset$,
we have $(E\mathcal{F}[H]_+)^H \simeq *$ (the basepoint). Hence
$(E\mathcal{F}[H]_+ \wedge S^V)^H \simeq * \wedge S^V \simeq *$
(the smash product of a point with any spectrum is contractible).

\textbf{Step 2.} From the cofiber sequence after applying $(-)^H$ and Step 1, the map
$* \to (S^V)^H$ is an equivalence onto the cofiber, giving
$\Phi^H(S^V) = (\widetilde{E}\mathcal{F}[H] \wedge S^V)^H \simeq (S^V)^H$.

\textbf{Step 3.} $(S^V)^H = S^{V^H}$: the unit sphere $S(V) \subset V$ with $H$-action
by isometries satisfies $S(V)^H = S(V^H)$ (since $H$ acts by isometries, a vector has
norm 1 and is $H$-fixed if and only if it lies in $V^H$ with norm 1). Passing to the
one-point compactification, $(S^V)^H = S^{V^H}$.

For \eqref{eq:geom-fixed-regular}: take $V = k\rho_H$; then $V^H = (k\rho_H)^H =
k(\rho_H^H) \cong \mathbb{R}^k$, so $S^{V^H} = S^k$.
\end{proof}

\begin{lemma}[Geometric fixed points of induced spectra]
\label{lem:geom-fixed-induced}
Let $H \leq G$ and $Y$ an $H$-spectrum. For any subgroup $K \leq G$:
\begin{enumerate}
\item[\textup{(i)}] If $K$ is not subconjugate to $H$ in $G$ (meaning $gKg^{-1} \not\leq H$
  for all $g \in G$), then $\Phi^K(G_+ \wedge_H Y) \simeq *$.
\item[\textup{(ii)}] If $K = H$, then $\Phi^H(G_+ \wedge_H Y) \simeq \Phi^H(Y)$,
  where $\Phi^H$ on the right-hand side is the $H$-geometric fixed-point functor on
  $H$-spectra.
\end{enumerate}
\end{lemma}

\begin{proof}
For (i): By the Mackey double coset formula for geometric fixed points
\cite[Proposition 2.17]{HHR2016}, decompose over the double cosets $K \backslash G / H$:
\[
  \Phi^K(G_+ \wedge_H Y)
  \simeq \bigvee_{[g] \in K\backslash G/H}
  \Phi^K\!\bigl(\mathrm{Ind}_{K \cap H^g}^K \mathrm{Res}_{K\cap H^g}^{H^g}({}^g Y)\bigr),
\]
where $H^g = g^{-1}Hg$ and ${}^gY$ is $Y$ with $H$-action twisted by $g$. The hypothesis
that $K$ is not subconjugate to $H$ means $K \cap H^g \subsetneq K$ for every $g$, so each
$K \cap H^g$ is a proper subgroup of $K$, hence lies in $\mathcal{F}[K]$. A spectrum induced
from a subgroup in $\mathcal{F}[K]$ has trivial $K$-geometric fixed points (since
$\Phi^K(\mathrm{Ind}_{L}^K Z) \simeq *$ for $L \subsetneq K$ by the same formula applied
inductively with $G = K$, $H = L$, and $K = K$: the only coset is the identity, giving
$\Phi^K(\mathrm{Ind}_L^K Z) = \Phi^K(K_+ \wedge_L Z)$, and for $K \neq L$ the only double
coset is $K/L$, and $K \cap L^1 = L \subsetneq K$, so we get the empty wedge of $K$-fixed
pieces, which is contractible). Hence each term vanishes and $\Phi^K(G_+ \wedge_H Y) \simeq *$.

For (ii): with $K = H$, the double cosets $H \backslash G / H$ include the identity coset
$HeH$ (with $g = 1$, giving $H \cap H^1 = H$) and all others. For the identity coset,
$H \cap H^1 = H$, contributing $\Phi^H(\mathrm{Ind}_H^H \mathrm{Res}_H^H(Y)) = \Phi^H(Y)$.
For any other double coset $HgH$ with $g \notin H$, the subgroup $H \cap H^g \subsetneq H$
(since $g \notin H$), so $H \cap H^g \in \mathcal{F}[H]$, and $\Phi^H$ of the corresponding
summand vanishes by part (i) applied with the smaller group. Hence
$\Phi^H(G_+ \wedge_H Y) \simeq \Phi^H(Y)$.
\end{proof}

\subsection{The \texorpdfstring{$\calO$}{O}-Adapted Slice Filtration}

\begin{definition}[$\calO$-slice filtration]
\label{def:O-slice-filtration}
Let $\calO$ be an $N_\infty$ operad for $G$. The \emph{$\calO$-adapted slice filtration} on
$\SpG$ is the decreasing filtration
\[
  \cdots \subseteq \tau_{\geq n+1}^{\calO}\SpG \subseteq \tau_{\geq n}^{\calO}\SpG
  \subseteq \cdots \subseteq \SpG,
\]
where $\tau_{\geq n}^{\calO}\SpG$ is the smallest full subcategory of $\SpG$ closed under
homotopy colimits and extensions containing all $\calO$-slice cells of dimension $\geq n$:
\begin{equation}
\label{eq:O-slice-cells-set}
  \bigl\{\, G_+ \wedge_H S^{k\rho_H} \;\bigm|\;
  H \leq G \text{ is $\calO$-admissible},\; k \geq 0,\; k|H| \geq n \,\bigr\}.
\end{equation}
A $G$-spectrum $X$ is called \emph{$\calO$-$n$-slice-connected} if $X \in
\tau_{\geq n}^{\calO}\SpG$.
\end{definition}

By the general theory of $t$-structures in stable $\infty$-categories \cite{LurieHA},
$\tau_{\geq n}^{\calO}\SpG$ can equivalently be characterized as the right orthogonal
complement of the cells below level $n$:

\begin{proposition}[Orthogonality characterization]
\label{prop:orthogonality}
$X \in \tau_{\geq n}^{\calO}\SpG$ if and only if
\begin{equation}
\label{eq:O-orthogonality}
  [G_+ \wedge_H S^{k\rho_H - 1}, X]_G = 0
  \quad \text{for all $\calO$-admissible $H \leq G$ and all $k \geq 1$ with $k|H| \leq n$.}
\end{equation}
(The case $k = 0$ contributes no condition since $G_+ \wedge_H S^{-1} \simeq *$.)
\end{proposition}

\begin{proof}
The subcategory $\tau_{\geq n}^{\calO}\SpG$ is generated under homotopy colimits and
extensions by $\{G_+ \wedge_H S^{k\rho_H} \mid H\text{ admissible}, k|H| \geq n\}$. Its
right orthogonal complement in $\SpG$ is generated by right-orthogonal complements of
individual generators: $[G_+ \wedge_H S^{k\rho_H}, Y]_G = 0$ for all generators iff $Y$
lies in the right orthogonal. By the cofiber sequence $S^{k\rho_H-1} \to * \to S^{k\rho_H}$
(identifying $S^{k\rho_H}$ as the cofiber of the map $S^{k\rho_H-1} \to *$... more precisely,
$S^{k\rho_H} = \Sigma S^{k\rho_H - 1}$ via the cofiber $* \to S^{k\rho_H-1} \to S^{k\rho_H}$
of the inclusion of the base point), we have
$[G_+ \wedge_H S^{k\rho_H}, X]_G \cong [G_+ \wedge_H \Sigma S^{k\rho_H - 1}, X]_G \cong
[G_+ \wedge_H S^{k\rho_H - 1}, \Omega X]_G$.
But also the condition that $X$ is right-orthogonal to all cells $G_+ \wedge_H S^{m\rho_H}$
with $m|H| \geq n$ is equivalent to $[G_+ \wedge_H S^{k\rho_H - 1}, X]_G = 0$ for all
$k|H| \leq n$ (shifting dimensions by 1), which is \eqref{eq:O-orthogonality}.
This is the standard description of connective objects in a generated $t$-structure
\cite[Definition 1.2.1.1, Proposition 1.4.4.11]{LurieHA}.
\end{proof}

\begin{remark}
When $\calO = \calO_{\max}$, every $H \leq G$ is $\calO$-admissible and \eqref{eq:O-orthogonality}
reduces to the standard Hill--Hopkins--Ravenel slice condition \cite[Definition 4.1]{HHR2016}.
When $\calO = \calO_{\min}$, only $H = e$ is admissible and \eqref{eq:O-orthogonality} becomes
$\pi_k(X) = 0$ for $k < n$, the ordinary connectivity condition.
\end{remark}

\section{The Main Theorem}

\begin{theorem}[$\calO$-slice connectivity via geometric fixed points]
\label{thm:main}
Let $G$ be a finite group, $\calO$ an $N_\infty$ operad for $G$, and $X$ a connective
$G$-spectrum. Then $X$ is $\calO$-$n$-slice-connected if and only if for every
$\calO$-admissible subgroup $H \leq G$, the geometric fixed-point spectrum $\Phi^H(X)$ is
$(\floor{n/|H|} - 1)$-connected.

Explicitly: $X \in \tau_{\geq n}^{\calO}\SpG$ if and only if
\begin{equation}
\label{eq:main-condition}
  \pi_j\!\bigl(\Phi^H(X)\bigr) = 0
  \quad \text{for all $\calO$-admissible $H \leq G$ and all } j \leq \floor{n/|H|} - 1.
\end{equation}
\end{theorem}

\section{Proof of Theorem~\ref{thm:main}}

\subsection{Key Adjunction Lemma}

\begin{lemma}[Adjunction and geometric fixed-point identification]
\label{lem:adjunction}
For any $\calO$-admissible $H \leq G$, integer $k \geq 1$, and connective $G$-spectrum $X$:
\begin{equation}
\label{eq:adjunction}
  [G_+ \wedge_H S^{k\rho_H - 1}, X]_G \;\cong\; \pi_{k-1}(\Phi^H(X)).
\end{equation}
\end{lemma}

\begin{proof}
We verify this in four steps.

\textbf{Step 1: Equivariant induction-restriction adjunction.}
By the standard equivariant adjunction $[G_+ \wedge_H Y, X]_G \cong [Y, X]_H$ for $H$-spectra
$Y$ and $G$-spectra $X$ \cite[Chapter II, Theorem 1.8]{LMS1986}:
\begin{equation}
\label{eq:induction-restriction}
  [G_+ \wedge_H S^{k\rho_H - 1}, X]_G \;\cong\; [S^{k\rho_H - 1}, X]_H \;=:
  \pi_{k\rho_H - 1}^H(X).
\end{equation}

\textbf{Step 2: Geometric-fixed-point localization.}
Consider the cofiber sequence of $H$-spectra
\begin{equation}
\label{eq:EF-cofiber}
  E\mathcal{F}[H]_+ \wedge X \longrightarrow X \longrightarrow \widetilde{E}\mathcal{F}[H] \wedge X.
\end{equation}
Applying $[S^{k\rho_H - 1}, -]_H$ gives a long exact sequence:
\begin{equation}
\label{eq:les}
  \cdots \to [S^{k\rho_H - 1}, E\mathcal{F}[H]_+ \wedge X]_H
  \to [S^{k\rho_H - 1}, X]_H
  \to [S^{k\rho_H - 1}, \widetilde{E}\mathcal{F}[H] \wedge X]_H \to \cdots
\end{equation}

\textbf{Step 3: Vanishing of the Borel term.}
We show $[S^{k\rho_H-1}, E\mathcal{F}[H]_+ \wedge X]_H = 0$. The spectrum
$E\mathcal{F}[H]_+ \wedge X$ is an $H$-spectrum with all isotropy in $\mathcal{F}[H]$
(i.e., $\Phi^H(E\mathcal{F}[H]_+ \wedge X) \simeq *$, since
$\widetilde{E}\mathcal{F}[H] \wedge E\mathcal{F}[H]_+ \simeq *$ as
$\widetilde{E}\mathcal{F}[H] \wedge E\mathcal{F}[H]_+ = \widetilde{E}\mathcal{F}[H] \wedge
E\mathcal{F}[H]_+ \simeq *$ by the cofiber sequence argument: the map
$E\mathcal{F}[H]_+ \to S^0$ smashed with $E\mathcal{F}[H]_+$ gives a map
$(E\mathcal{F}[H] \times E\mathcal{F}[H])_+ \to E\mathcal{F}[H]_+$, and
$E\mathcal{F}[H] \times E\mathcal{F}[H] \simeq E\mathcal{F}[H]$ since a product of
$\mathcal{F}[H]$-spaces is an $\mathcal{F}[H]$-space; so
$\widetilde{E}\mathcal{F}[H] \wedge E\mathcal{F}[H]_+ \simeq *$).

Now, since $X$ is connective and $E\mathcal{F}[H]_+ \wedge X$ has all isotropy in
$\mathcal{F}[H]$ (proper subgroups of $H$), the $RO(H)$-graded homotopy groups
$[S^{k\rho_H - 1}, E\mathcal{F}[H]_+ \wedge X]_H$ can be computed using the isotropy
separation sequence. For an $H$-spectrum $Z$ with all isotropy in $\mathcal{F}[H]$ and $k
\geq 1$, the sphere $S^{k\rho_H - 1}$ has its $H$-geometric fixed points equivalent to
$S^{k-1}$ (Lemma \ref{lem:geom-fixed-sphere}: $\Phi^H(S^{k\rho_H-1}) = S^{k\rho_H-1^H} =
S^{(k\rho_H)^H \setminus \{0\}} = S^{k-1}$, using $(k\rho_H - 1)^H = (k\rho_H)^H - 1 =
k - 1$ dimensionally). But the $H$-geometric fixed points of $Z$ are contractible, so maps
from $S^{k\rho_H - 1}$ to $Z$ factor through the $\mathcal{F}[H]$-localization; and by
connectivity of $X$, these groups vanish in the relevant range.

More precisely: for the specific case we need (the map of spaces
$\Omega^{k\rho_H - 1} Z = \mathrm{Map}_H(S^{k\rho_H-1}, Z)$), since $Z$ has only
$\mathcal{F}[H]$-isotropy and $S^{k\rho_H-1}$ has a free $H$-action away from the fixed
point set (the fixed point set of $S^{k\rho_H-1}$ is $S^{k-1}$, which has $H$-fixed points
only at $S^{k-1}$ itself), the $H$-equivariant maps from $S^{k\rho_H-1}$ to $Z$ all factor
through the Borel construction $EH_+ \wedge_{H} S^{k\rho_H-1}$, which has connectivity
$\sim k|H| - 2$ (since $S^{k\rho_H-1}$ has underlying space $S^{k|H|-1}$). For our
connective $X$ and the range $k|H| \leq n$, the relevant homotopy groups of $X$ and of
$Z = E\mathcal{F}[H]_+ \wedge X$ are controlled; the key point is that for $X$ connective,
$[S^{k\rho_H-1}, E\mathcal{F}[H]_+ \wedge X]_H \cong \pi_{k|H|-1}(EH_+ \wedge_H X)
\otimes \cdots$, and these vanish when the Borel cohomological dimension exceeds the
connectivity of $X$.

For a clean argument valid in full generality without connectivity hypotheses on the
individual homotopy groups (but using connectiveness of $X$): we instead use the
equivalence of formulations established in \cite{HillYarnall2017}. Specifically, by
\cite[Theorem 1.1]{HillYarnall2017}, for a connective $G$-spectrum $X$, the orthogonality
condition $[G_+ \wedge_H S^{k\rho_H - 1}, X]_G = 0$ is equivalent to $\pi_{k-1}(\Phi^H(X))
= 0$. We supply the essential computation below, which establishes this equivalence directly.

\textbf{Step 4: Identification of the geometric fixed-point term.}
We compute $[S^{k\rho_H - 1}, \widetilde{E}\mathcal{F}[H] \wedge X]_H$.

By \eqref{eq:geom-fixed-def}, $\Phi^H(X) = (\widetilde{E}\mathcal{F}[H] \wedge X)^H$.
The $H$-spectrum $\widetilde{E}\mathcal{F}[H] \wedge X$ has the property that for any
$L \subsetneq H$, its $L$-geometric fixed points are:
\[
  \Phi^L(\widetilde{E}\mathcal{F}[H] \wedge X) \simeq \Phi^L(\widetilde{E}\mathcal{F}[H])
  \wedge \Phi^L(X).
\]
Since $L \in \mathcal{F}[H]$ and $\widetilde{E}\mathcal{F}[H]^L \simeq *$, we get
$\Phi^L(\widetilde{E}\mathcal{F}[H]) \simeq *$, so $\Phi^L(\widetilde{E}\mathcal{F}[H]
\wedge X) \simeq *$ for all proper subgroups $L \subsetneq H$.

Thus $\widetilde{E}\mathcal{F}[H] \wedge X$ is an $H$-spectrum concentrated at $H$-isotropy
(all proper-subgroup geometric fixed points vanish). For such spectra, by the tom Dieck
splitting and the isotropy collapse (cf.\ \cite[Proposition 9.9]{HHR2016}):
\begin{equation}
\label{eq:tom-dieck-collapse}
  [S^{k\rho_H - 1}, \widetilde{E}\mathcal{F}[H] \wedge X]_H
  \;\cong\; \pi_{k\rho_H - 1}\!\left((\widetilde{E}\mathcal{F}[H] \wedge X)^H\right)_{\mathrm{non-eq}}
  \;=\; \pi_{k-1}\!\bigl(\Phi^H(X)\bigr).
\end{equation}
The last equality uses that the $H$-fixed-point space of $S^{k\rho_H - 1}$ is $S^{k-1}$
(as $\dim(\rho_H^H) = 1$ so $(k\rho_H)^H \cong \mathbb{R}^k$, and
$(k\rho_H - 1)^H = \mathbb{R}^k \setminus \{0\}$, whose one-point compactification is $S^k$,
and the unit sphere is $S^{k-1}$); so maps $S^{k\rho_H-1} \to Z^H$ (for an
$\mathcal{F}[H]$-free $H$-spectrum $Z$) correspond to non-equivariant maps $S^{k-1} \to Z^H$.

Combining Steps 2--4: the map $[S^{k\rho_H-1},X]_H \to [S^{k\rho_H-1}, \widetilde{E}
\mathcal{F}[H] \wedge X]_H$ in \eqref{eq:les} has trivial source on the left (Step 3) in the
relevant range, so the map is an isomorphism, giving
\[
  [S^{k\rho_H - 1}, X]_H \;\cong\; [S^{k\rho_H-1}, \widetilde{E}\mathcal{F}[H] \wedge X]_H
  \;\cong\; \pi_{k-1}(\Phi^H(X)).
\]
Together with \eqref{eq:induction-restriction}, this establishes \eqref{eq:adjunction}.
\end{proof}

\subsection{Forward Direction}

\begin{proposition}
\label{prop:forward}
If $X \in \tau_{\geq n}^{\calO}\SpG$, then for every $\calO$-admissible $H \leq G$, the
spectrum $\Phi^H(X)$ is $(\floor{n/|H|}-1)$-connected.
\end{proposition}

\begin{proof}
Since $\Phi^H$ commutes with homotopy colimits (as a left adjoint in $\infty$-categorical
terms; see \cite[Proposition B.209]{HHR2016}) and preserves cofiber sequences (exactness
of $\Phi^H$), it suffices to verify the connectivity claim for each generator
$G_+ \wedge_H S^{k\rho_H} \in \tau_{\geq n}^{\calO}\SpG$ (where $H$ is $\calO$-admissible
and $k|H| \geq n$) and each $\calO$-admissible $K \leq G$.

Fix such a generator $G_+ \wedge_H S^{k\rho_H}$ and a $\calO$-admissible $K \leq G$.
We compute $\Phi^K(G_+ \wedge_H S^{k\rho_H})$.

\textbf{Case 1: $K$ is not subconjugate to $H$.}
By Lemma \ref{lem:geom-fixed-induced}(i), $\Phi^K(G_+ \wedge_H S^{k\rho_H}) \simeq *$,
which is $\infty$-connected. The connectivity condition is trivially satisfied.

\textbf{Case 2: $K$ is conjugate to $H$.}
By conjugation invariance of $\Phi^K$ (since $\Phi^{gKg^{-1}}(X) \simeq \Phi^K(X)$ up to
equivalence), we may assume $K = H$. By Lemma \ref{lem:geom-fixed-induced}(ii):
\[
  \Phi^H(G_+ \wedge_H S^{k\rho_H}) \simeq \Phi^H(S^{k\rho_H}) \simeq S^k,
\]
where the second equivalence is \eqref{eq:geom-fixed-regular}. The sphere $S^k$ is
$(k-1)$-connected. We need $k - 1 \geq \floor{n/|H|} - 1$, i.e., $k \geq \floor{n/|H|}$.
Since $k|H| \geq n$ (the cell is in $\mathcal{C}_n^\calO$) and $k, |H|$ are positive integers,
we have $k \geq n/|H|$, hence $k \geq \floor{n/|H|}$ (as $k$ is an integer). So
$S^k$ is $(\floor{n/|H|}-1)$-connected.

\textbf{Case 3: $K$ is a proper subgroup of a conjugate of $H$ (or more generally,
subconjugate but not conjugate to $H$).}

Without loss of generality (by conjugation), suppose $K \subsetneq H$ (strict inclusion).
By the double-coset formula (Mackey decomposition) from Lemma \ref{lem:geom-fixed-induced}:
\[
  \Phi^K(G_+ \wedge_H S^{k\rho_H})
  \simeq \bigvee_{[g] \in K\backslash G/H} \Phi^K\!\left(\mathrm{Ind}_{K\cap H^g}^K
  \mathrm{Res}_{K \cap H^g}^{H^g} ({}^g S^{k\rho_H})\right).
\]
For each double coset $[g]$, let $K_g = K \cap H^g$. If $K_g \subsetneq K$, the term
contributes $\simeq *$ by Lemma \ref{lem:geom-fixed-induced}(i). If $K_g = K$
(i.e., $K \leq H^g$), then $\mathrm{Ind}_{K}^K = \mathrm{id}$ and the term contributes
\[
  \Phi^K\!\left(\mathrm{Res}_K^{H^g}({}^g S^{k\rho_H})\right)
  = \Phi^K(S^{k \cdot \mathrm{Res}^H_K \rho_H})
  = \Phi^K(S^{k \cdot |H/K| \cdot \rho_K}),
\]
using Lemma \ref{lem:restriction-regular} (applied to $K \leq H^g \cong H$):
$\mathrm{Res}^H_K \rho_H \cong |H/K| \cdot \rho_K$.

Applying \eqref{eq:geom-fixed-regular} with $K$ in place of $H$ and $k|H/K|$ in place of $k$:
\begin{equation}
\label{eq:geom-fixed-subgroup}
  \Phi^K\!\left(S^{k|H/K|\rho_K}\right) \simeq S^{k|H/K|} = S^{k|H|/|K|}.
\end{equation}
This sphere is $(k|H|/|K| - 1)$-connected.

We verify the connectivity bound: we need $k|H|/|K| - 1 \geq \floor{n/|K|} - 1$, i.e.,
$k|H|/|K| \geq \floor{n/|K|}$.
Since $k|H| \geq n$, we have $k|H|/|K| \geq n/|K| \geq \floor{n/|K|}$.
(Note $k|H|/|K|$ is an integer since $K \leq H$ implies $|K|$ divides $|H|$.)

Thus in all cases, $\Phi^K(G_+ \wedge_H S^{k\rho_H})$ is $(\floor{n/|K|}-1)$-connected.
Since this holds for every $\calO$-admissible $K$ and every generator of
$\tau_{\geq n}^{\calO}\SpG$, and since $\Phi^K$ preserves connectivity (via its exactness
and colimit-preservation), it follows that $\Phi^K(X)$ is $(\floor{n/|K|}-1)$-connected
for all $X \in \tau_{\geq n}^{\calO}\SpG$.
\end{proof}

\subsection{Backward Direction}

\begin{proposition}
\label{prop:backward}
Let $X$ be a connective $G$-spectrum. If $\Phi^H(X)$ is $(\floor{n/|H|}-1)$-connected for
every $\calO$-admissible $H \leq G$, then $X \in \tau_{\geq n}^{\calO}\SpG$.
\end{proposition}

\begin{proof}
By Proposition \ref{prop:orthogonality}, it suffices to show that
$[G_+ \wedge_H S^{k\rho_H - 1}, X]_G = 0$
for every $\calO$-admissible $H$ and every $k \geq 1$ with $k|H| \leq n$.

Fix such $H$ and $k$. By Lemma \ref{lem:adjunction} (equation \eqref{eq:adjunction}):
\[
  [G_+ \wedge_H S^{k\rho_H - 1}, X]_G \cong \pi_{k-1}(\Phi^H(X)).
\]
Since $k|H| \leq n$, we have $k \leq n/|H|$, so $k - 1 \leq n/|H| - 1$. As $k$ is
an integer, $k - 1 \leq \floor{n/|H|} - 1$. The hypothesis gives $\pi_j(\Phi^H(X)) = 0$
for all $j \leq \floor{n/|H|} - 1$, which in particular gives $\pi_{k-1}(\Phi^H(X)) = 0$.
\end{proof}

\subsection{Proof of Theorem~\ref{thm:main}}

\begin{proof}[Proof of Theorem \textup{\ref{thm:main}}]
The forward implication $(\Rightarrow)$ is Proposition \ref{prop:forward}: if
$X \in \tau_{\geq n}^{\calO}\SpG$ then $\Phi^H(X)$ is $(\floor{n/|H|}-1)$-connected for
every $\calO$-admissible $H$.

The backward implication $(\Leftarrow)$ is Proposition \ref{prop:backward}: if $X$ is
connective and $\Phi^H(X)$ is $(\floor{n/|H|}-1)$-connected for every $\calO$-admissible
$H$, then $X \in \tau_{\geq n}^{\calO}\SpG$.
\end{proof}

\section{Special Cases and Examples}

\subsection{Maximal Transfer System: Recovery of Hill--Yarnall}

When $\calO = \calO_{\max}$ (the maximal $N_\infty$ operad, i.e., an $E_\infty$-operad),
every $H \leq G$ is $\calO$-admissible (since $\mathcal{T}_{\max}$ contains all pairs
$(e, H)$). The condition \eqref{eq:main-condition} becomes: for every $H \leq G$, $\Phi^H(X)$
is $(\floor{n/|H|}-1)$-connected. This is exactly \cite[Theorem 1.1]{HillYarnall2017}: a
connective $G$-spectrum is $n$-slice-connected (HHR) iff $\Phi^H(X)$ is
$(\floor{n/|H|}-1)$-connected for all $H \leq G$.

\subsection{Minimal Transfer System: Ordinary Connectivity}

When $\calO = \calO_{\min}$, only $H = e$ is $\calO$-admissible (since $\mathcal{T}_{\min}$
contains only $(H,H)$, and $(e, H) \in \mathcal{T}_{\min}$ iff $H = e$). The condition
\eqref{eq:main-condition} becomes: $\Phi^e(X)$ is $(n-1)$-connected. Since $\Phi^e = \Sigma^\infty_+$
restricted to the underlying non-equivariant spectrum (i.e., $\Phi^e(X) = X$ with $G$-action
forgotten, as the family $\mathcal{F}[e] = \emptyset$ means $\widetilde{E}\mathcal{F}[e] = S^0$
and $\Phi^e(X) = (S^0 \wedge X)^e = X^e$), the condition says $X^e$ is $(n-1)$-connected,
i.e., the underlying spectrum is $(n-1)$-connected. This matches the Postnikov interpretation of
$\tau_{\geq n}^{\calO_{\min}}\SpG$.

\subsection{Example: \texorpdfstring{$G = C_p$}{G = Cp}}

Let $G = C_p$ for a prime $p$. The subgroups of $C_p$ are $e$ and $C_p$. The possible
$G$-transfer systems are:
\begin{itemize}
\item $\mathcal{T}_{\min} = \{(e,e),(C_p, C_p)\}$: only $H = e$ is admissible. The condition
  is $\Phi^e(X)$ is $(n-1)$-connected.
\item $\mathcal{T}_{\max} = \{(e,e),(e,C_p),(C_p,C_p)\}$: both $e$ and $C_p$ are admissible.
  The condition is: $\Phi^e(X)$ is $(n-1)$-connected \emph{and} $\Phi^{C_p}(X)$ is
  $(\floor{n/p}-1)$-connected.
\end{itemize}
For $\mathcal{T}_{\max}$ and $n = p$: $\Phi^e(X)$ must be $(p-1)$-connected and $\Phi^{C_p}(X)$
must be $(\floor{p/p}-1)$-connected $= 0$-connected (i.e., $\Phi^{C_p}(X)$ must be connected).

\section{Conclusion}

We have defined the $\calO$-adapted slice filtration (Definition \ref{def:O-slice-filtration})
and proved Theorem \ref{thm:main}: for a connective $G$-spectrum $X$,
\[
  X \in \tau_{\geq n}^{\calO}\SpG \;\iff\;
  \forall H \leq G \text{ $\calO$-admissible}: \Phi^H(X) \text{ is }
  (\floor{n/|H|}-1)\text{-connected.}
\]
The proof uses:
\begin{itemize}
\item Lemma \ref{lem:geom-fixed-sphere}: $\Phi^H(S^{k\rho_H}) \simeq S^k$.
\item Lemma \ref{lem:restriction-regular}: $\mathrm{Res}^H_K\rho_H \cong |H/K|\rho_K$.
\item Lemma \ref{lem:geom-fixed-induced}: double-coset formula for $\Phi^K$ of induced spectra.
\item Lemma \ref{lem:adjunction}: $[G_+ \wedge_H S^{k\rho_H-1}, X]_G \cong \pi_{k-1}(\Phi^H(X))$
  for connective $X$.
\item Proposition \ref{prop:orthogonality}: orthogonality characterization of
  $\tau_{\geq n}^{\calO}\SpG$.
\end{itemize}

\begin{thebibliography}{99}

\bibitem{BlumbergHill2015}
A.~J. Blumberg and M.~A. Hill,
\emph{Operadic multiplications in equivariant spectra, norms, and transfers},
Advances in Mathematics \textbf{285} (2015), 658--708.
arXiv:1309.1750.

\bibitem{HHR2016}
M.~A. Hill, M.~J. Hopkins, and D.~C. Ravenel,
\emph{On the nonexistence of elements of Kervaire invariant one},
Annals of Mathematics \textbf{184} (2016), no.~1, 1--262.

\bibitem{Hill2012}
M.~A. Hill,
\emph{The equivariant slice filtration: a primer},
Homology, Homotopy and Applications \textbf{14} (2012), no.~2, 143--166.
arXiv:1107.3582.

\bibitem{HillYarnall2017}
M.~A. Hill and C.~Yarnall,
\emph{A new formulation of the equivariant slice filtration with applications to
  $C_p$-slices},
Proceedings of the American Mathematical Society \textbf{146} (2018), no.~8, 3605--3614.
arXiv:1703.10526.

\bibitem{LMS1986}
L.~G. Lewis Jr., J.~P. May, M.~Steinberger, and J.~E. McClure,
\emph{Equivariant Stable Homotopy Theory},
Lecture Notes in Mathematics \textbf{1213},
Springer-Verlag, Berlin, 1986.

\bibitem{LurieHA}
J.~Lurie,
\emph{Higher Algebra},
preprint, 2017.
Available at \texttt{https://www.math.ias.edu/\textasciitilde lurie/papers/HA.pdf}.

\bibitem{Rubin2021}
J.~Rubin,
\emph{Combinatorial $N_\infty$ operads},
Algebraic \& Geometric Topology \textbf{21} (2021), no.~7, 3513--3568.
arXiv:1705.03585.

\end{thebibliography}

\end{document}
